\chapter{Lecture 6: File and String Streams, Overloading, References}
\begin{center}
\textit{January 23, 2018}
\end{center}

\section{File Streams}
The stream abstraction can be used on other sources of data:
\begin{lstlisting}[style=cpp]
#include <fstream>
// std::ifstream -- read from a file
// std::ofstream -- write to a file
\end{lstlisting}

\begin{lstlisting}[style=cpp]
// io/readFile.cc
#include <iostream>
#include <fstream>
#include <string>
using namespace std;

int main() {
    ifstream filestream{"myfile.txt"};   // "myfile.txt" is a path
    string s;
    while (filestream >> s) {
        cout << s << endl;
    }
}
\end{lstlisting}
The file is opened when \code{filestream} is initialized.

\paragraph{Q.} When is it closed?

\paragraph{A.} We'll look at this in detail about five weeks from now: it is
closed when the \code{filestream} variable goes out of scope (here, when
\code{main} returns). No explicit \code{close} needed.

Constructor parameters exist for appending, binary mode and so on. The key point:

\begin{important}[Streams are interchangeable]
Anything we can do with an \code{istream} (\code{ostream}) variable such as
\code{cin} (\code{cout}), we can do with an \code{ifstream} (\code{ofstream})
variable.
\end{important}

\section{String Streams}
We can also attach streams to strings, to read from or write to them:
\begin{lstlisting}[style=cpp]
#include <sstream>
// std::istringstream -- read from a string
// std::ostringstream -- write to a string
\end{lstlisting}

Building a string with output operations:
\begin{lstlisting}[style=cpp]
int low = 1, high = 100;
ostringstream ss;
ss << "Enter a value between " << low << " and " << high;
string s = ss.str();     // extract the accumulated string
\end{lstlisting}

Converting a string to a number, safely:
\begin{lstlisting}[style=cpp]
// io/getNum.cc
#include <iostream>
#include <sstream>
#include <string>
using namespace std;

int main() {
    int n;
    while (true) {
        cout << "Enter a number" << endl;
        string s;
        cin >> s;                 // always succeeds (other than at EOF)
        istringstream ss{s};
        if (ss >> n) break;
    }
    cout << "You entered " << n << endl;
}
\end{lstlisting}
Benefits:
\begin{itemize}
  \item Anything the user types can be read into a string, so the read from
        \code{cin} only fails at EOF. (Exercise: as written, this loops forever
        if the user gives up and sends EOF; add a check.)
  \item If the read from \code{ss} fails, we don't need \code{clear} and
        \code{ignore}: we just make a fresh \code{istringstream} next time around.
\end{itemize}
This is often the preferred way to parse input.

\begin{lstlisting}[style=cpp]
// readIntsSS.cc: echo all ints, skipping non-ints
int main() {
    string s;
    while (cin >> s) {
        istringstream ss{s};
        int n;
        if (ss >> n) cout << n << endl;
    }
}
\end{lstlisting}

\begin{aside}
\textbf{Homework (not for marks):} compare \code{readInts5.cc} (the
\code{clear}/\code{ignore} version from last lecture) with \code{readIntsSS.cc}
on the input \code{hello123abc45}. The first skips one character at a time, so it
eventually reads \code{123} and \code{45}. The second reads the whole word
\code{hello123abc45} as one string; \code{ss >> n} fails at \code{h}, so
nothing is printed. (On \code{123abc} the second version prints \code{123}.)
\end{aside}

\section{Default Arguments}
\begin{lstlisting}[style=cpp]
void printFile(string name = "myfile.txt") {
    ifstream filestream{name};
    string s;
    while (filestream >> s) cout << s << endl;
}

printFile("other.txt");   // prints other.txt
printFile();              // uses the default argument: myfile.txt
\end{lstlisting}

Default arguments must always appear last:
\begin{lstlisting}[style=cpp]
void test(int x = 5, string str) {}          // error: won't compile
void foo(int x = 5, string str = "Hello") {} // OK
\end{lstlisting}
When you call a function and omit arguments, you can only omit the last one (or
the last $n$):
\begin{lstlisting}[style=cpp]
foo(10, "World");   // OK
foo(10);            // OK, str = "Hello"
foo();              // OK, x = 5, str = "Hello"
foo("World");       // error: can't skip x
\end{lstlisting}
The default is filled in by the \emph{caller} at compile time, which is why it
belongs in the declaration (the header), not the definition.

\section{Function Overloading}
\paragraph{Q.} Can two functions have the same name?

\paragraph{A.} Not in C, but yes in C++. This is \textbf{function overloading}.

C++ allows functions with the same name as long as the number or types of
parameters differ. Differing only in return type is not enough.
\begin{lstlisting}[style=cpp]
int neg(int x)   { return -x; }
bool neg(bool b) { return !b; }
\end{lstlisting}
The compiler picks the version based on the argument types, at compile time.

\paragraph{Q.} Given \code{foo(int x = 5, string str = "Hello")} above, could we
also implement \code{void foo(int x)}?

\paragraph{A.} Declaring it is legal, but then the call \code{foo(10)} is
\emph{ambiguous} (both versions match equally well), so it won't compile. Default
arguments and overloading must not overlap.

\subsection{Operators Are Functions}
\begin{lstlisting}[style=cpp]
int a = 21, b = 3;
a >> b;    // legal: bit shift
cin >> x;  // legal: read
\end{lstlisting}
C++ implements operators as functions. The expression \code{cin >> x} is a call
to a function named \code{operator>>}, overloaded on the types of its operands.

\section{Pointers and const}
\begin{quote}
\textit{``Now we have a big NOTE in my NOTES.''} --- Nomair Naeem
\end{quote}
Read declarations from right to left:
\begin{lstlisting}[style=cpp]
int n = 5, x = 7;

const int *p = &n;     // p is a pointer to a const int
p = &x;                // fine: p itself can change
*p = 10;               // illegal: can't change the int through p
n = 10;                // fine: n itself isn't const

int * const q = &n;    // q is a const pointer to an int
q = &x;                // illegal: q can't be re-pointed
*q = 10;               // fine
\end{lstlisting}

\section{Structures in C++}
\begin{lstlisting}[style=cpp]
struct Node {
    int data;
    Node *next;        // no need to write "struct Node *next"
};                     // don't forget the semicolon!

Node n = {3, nullptr}; // no "struct" here either
\end{lstlisting}
Do not use \code{NULL} or \code{0} for pointers in C++: C++11 introduced the
keyword \code{nullptr}.
\begin{quote}
\textit{``Use \code{nullptr}, as I sometimes tend to go back to old habits.''}
--- Nomair Naeem
\end{quote}

\section{Parameter Passing}
\begin{lstlisting}[style=cpp]
void inc(int n) { n += 1; }

int x = 5;
inc(x);
cout << x << endl;   // prints 5
\end{lstlisting}
In C and C++, arguments are passed by value by default: \code{inc} gets a copy.
In C we would pass a pointer:
\begin{lstlisting}[style=cpp]
void inc(int *n) { *n += 1; }

int x = 5;
inc(&x);
cout << x << endl;   // prints 6
\end{lstlisting}
This is why C needs \code{scanf("\%d", \&x)}. But in C++ we write
\code{cin >> x}, not \code{cin >> \&x}. How does \code{cin} change \code{x}?

\section{lvalue References}
C++ has a \textbf{reference} type, written \code{\&}.
\begin{lstlisting}[style=cpp]
int x = 10;
int &z = x;          // z is an lvalue reference to x
int * const p = &x;  // z behaves much like this constant pointer...

z = 5;               // ... with automatic dereferencing: same as *p = 5
                     // (NOT *z = 5)
\end{lstlisting}

\begin{definition}
An \textbf{lvalue reference} is like a constant pointer with automatic
dereferencing. It is another name (an \emph{alias}) for the object it refers to.
\end{definition}

\begin{center}
\begin{tikzpicture}[box/.style={draw, minimum width=12mm, minimum height=7mm}]
\node[box] (x) {10};
\node[left=2mm of x] {\code{x}, \code{z}};
\node[box, right=25mm of x] (p) {};
\node[right=2mm of p] {\code{p}};
\draw[-{Stealth}] (p.center) -- (x);
\end{tikzpicture}
\end{center}

If you take the address of \code{z}, you get the address of \code{x}; if you ask
for \code{sizeof(z)} you get \code{sizeof(int)}. \code{z} has no identity of its
own; it is only a name for \code{x}. (This is called \emph{aliasing}.)

\subsection{Restrictions on References}
\begin{itemize}
  \item References must be initialized: \code{int \&z;} is an error.
  \item They must be initialized to an \textbf{lvalue}, i.e.\ something that
        denotes a storage location (has an address):
\begin{lstlisting}[style=cpp]
int &z = x;      // OK
int &w = y + y;  // error: y + y has no address
int &v = 5;      // error
\end{lstlisting}
  \item Once initialized, a reference can never be made to refer to something
        else.
  \item You cannot create a pointer to a reference (\code{int \&*p} is illegal;
        a reference to a pointer, \code{int *\&r = p;}, is fine).
  \item You cannot create a reference to a reference.
  \item You cannot create an array of references.
\end{itemize}

\subsection{Pass by Reference}
\begin{lstlisting}[style=cpp]
void inc(int &n) { n = n + 1; }   // no pointer syntax needed

int x = 5;
inc(x);
cout << x << endl;   // prints 6
\end{lstlisting}
\code{n} is another name for \code{x}, so changing \code{n} changes \code{x}.
This is how \code{cin >> x} works: \code{operator>>} takes \code{x} by
reference.

\section*{Summary}
\begin{itemize}
  \item \code{ifstream}/\code{ofstream} and \code{istringstream}/\code{ostringstream}
        work exactly like \code{cin}/\code{cout}; files close automatically when
        the stream goes out of scope.
  \item Reading a string then parsing it with an \code{istringstream} is a
        robust way to handle bad input.
  \item Default arguments go last; overloading requires different parameter
        lists; operators are functions named \code{operator<op>}.
  \item Read \code{const} declarations right to left.
  \item Use \code{nullptr}. A reference is an alias: a constant pointer with
        automatic dereferencing, must be initialized to an lvalue.
\end{itemize}
